% Revised version of "Q-EVM: A Quantum-Resistant Account Abstraction Layer for Ethereum".
% The original paper PDF is unchanged and kept outside this repository; see CHANGES.md for every change.
% Numbers come from tables/*.tex and data/*.csv, generated by scripts/make_paper_tables.py
% from results/ -- do not type numbers by hand.
\documentclass[conference]{IEEEtran}
\usepackage[T1]{fontenc}
\usepackage{amsmath,amssymb}
\usepackage{booktabs}
\usepackage{graphicx}
\usepackage{algorithm}
\usepackage{algpseudocode}
\usepackage{pgfplots}
\usepackage{xcolor}
\usepackage[hidelinks]{hyperref}
\pgfplotsset{compat=1.17}

\newcommand{\NativeModel}{2{,}802{,}020}
\newcommand{\SingleOpGas}{251{,}302}
\newcommand{\BatchEightGas}{39{,}387}
\newcommand{\BatchMaxGas}{10{,}060}
\newcommand{\BatchMaxN}{256}
\newcommand{\AttestMaxGas}{9{,}215}
\newcommand{\AttestOneGas}{34{,}953}
\newcommand{\PerOpGas}{251{,}297}
\newcommand{\GrothOverhead}{216{,}344}
\newcommand{\RedModelOne}{91.0}
\newcommand{\RedModelMax}{99.6}
\newcommand{\ProveMsOne}{24.1}
\newcommand{\ProveMsMax}{23.1}
\newcommand{\MldsaVerifyMsPerOp}{0.030}
\newcommand{\EtoEMsOne}{38.7}
\newcommand{\EtoEMsMax}{68.0}
\newcommand{\EtoEPerOpMax}{0.27}
\newcommand{\HeapMax}{10.3}
\newcommand{\Repeats}{30}
\newcommand{\NativeExecMain}{8{,}168{,}152}
\newcommand{\NativeExecMainRaw}{8168152}
\newcommand{\RedMainOne}{96.9}
\newcommand{\RedMainMax}{99.88}
\newcommand{\MainVsModel}{-65.7}
\newcommand{\NativeExec}{1{,}248{,}996}
\newcommand{\NativeExecRaw}{1248996}
\newcommand{\RedExecOne}{79.9}
\newcommand{\RedExecMax}{99.19}
\newcommand{\SetKeyGas}{4{,}837{,}624}
\newcommand{\ExecVsModel}{+124.3}
\newcommand{\SecCases}{45}
\newcommand{\SecPass}{45}
\newcommand{\ProverUsdPerOpProj}{0.0035}
\newcommand{\SavedUsdOneGwei}{3.72}
\newcommand{\BreakEvenTenthGwei}{1{,}330}
\newcommand{\BreakEvenOneGwei}{13{,}301}
\newcommand{\EthUsd}{3{,}000}
\newcommand{\ProverUsdHour}{1.006}

\providecommand{\NativeExec}{n/a}
\providecommand{\RedExecOne}{n/a}
\providecommand{\RedExecMax}{n/a}
\providecommand{\ExecVsModel}{n/a}
\newif\ifexecbaseline
\IfFileExists{tables/native_baseline.tex}{\execbaselinetrue}{\execbaselinefalse}

\begin{document}

\title{Q-EVM: Forkless Post-Quantum Account Abstraction for Ethereum --- An Experimental Study of the Gas, Batching, Latency and Trust Trade-off}

\author{\IEEEauthorblockN{Mohit Anuragi, Srikanth Reddy}
\IEEEauthorblockA{Department of Computer Science and Engineering, Indian Institute of Information Technology Kota, India\\
Email: \{2023kucp1155, 2023kucp1151\}@iiitkota.ac.in}
\and
\IEEEauthorblockN{Supervisor: Dr.\ Isha Pathak Tripathi}
\IEEEauthorblockA{Assistant Professor, Department of CSE, IIIT Kota, India}}

\maketitle

\begin{abstract}
Verifying post-quantum signatures such as ML-DSA (FIPS~204) directly in the Ethereum Virtual Machine costs on the order of a million gas or more per signature. We study a forkless alternative built on account abstraction. A bundler verifies ML-DSA-44 signatures off-chain, and an on-chain wallet contract accepts a batch of user operations against a single Groth16 receipt. The wallet binds keys to accounts, consumes nonces and recomputes every operation hash and the batch digest on-chain.

We are explicit about what is proven. In the implemented system the Groth16 circuit binds the batch digest but does \emph{not} prove ML-DSA verification, which the bundler performs and is trusted for. Proving ML-DSA inside a zkVM is left as future work, with a concrete plan.

We report repeated measurements ($n=\Repeats$ per point) on a real EVM implementation (revm, Cancun rules) across batch sizes $N=1$ to $\BatchMaxN$, together with an ablation that separates off-chain verification, Groth16 verification and batching, and a \SecCases-case security evaluation. As a baseline we execute open-source Solidity ML-DSA-44 verifiers on the same chain; they cost \NativeExec\ to \NativeExecMain\ gas per signature. A single-operation Q-EVM transaction costs \SingleOpGas\ gas (\RedExecOne\% less than the cheapest executed verifier). Batching reduces this to \BatchMaxGas\ gas per operation at $N=\BatchMaxN$ (\RedExecMax\% less), close to the per-operation floor (\AttestMaxGas\ gas) of a trusted-bundler attestation. A sound zkVM proof would cost far less than the gas it saves, so its proving \emph{latency}, not its price, is the binding constraint.
\end{abstract}

\begin{IEEEkeywords}
Post-quantum cryptography, ML-DSA, account abstraction, ERC-4337, Groth16, zero-knowledge proofs, Ethereum, gas cost.
\end{IEEEkeywords}

\section{Introduction}

\subsection{Motivation}
Ethereum accounts are authorised with ECDSA over secp256k1, whose security rests on the elliptic-curve discrete-logarithm problem. Shor's algorithm solves this problem in polynomial time on a sufficiently large quantum computer~\cite{shor1994}. NIST has standardised the lattice-based ML-DSA (formerly CRYSTALS-Dilithium) as FIPS~204~\cite{fips204}. Verifying ML-DSA in the EVM is expensive, however. The EVM has no SHAKE or NTT primitives, and open-source Solidity verifiers report about 1.2M gas~\cite{fireblocks2026evmmldsa,zknox2026ethdilithium} to 8.1M gas~\cite{zknox2026ethdilithium} per ML-DSA-44 verification (unaudited, 2026). Precompiles have been proposed (EIP-8051~\cite{eip8051}, EIP-8355~\cite{eip8355}), but they require a hard fork.

\subsection{Problem and approach}
We ask how much on-chain cost can be removed \emph{without a fork}, what that costs off-chain and in latency, and what security is retained. In Q-EVM, users sign ERC-4337-style~\cite{eip4337} user operations with ML-DSA-44. A bundler verifies them off-chain, and an on-chain wallet contract accepts a whole batch against one succinct receipt.

\subsection{Scope: what is implemented}
Table~\ref{tab:status} lists what is implemented, measured, modelled, projected and future. In particular, \textbf{ML-DSA verification runs off-chain in the bundler}. The Groth16 receipt verified on-chain binds the batch digest but does not prove that the signature was valid (Section~\ref{sec:proof}). An earlier draft of this work reported a RISC~Zero cycle count and a 12.4\,s proving time. Neither is reproducible from our code, so neither is presented here as a measurement.

\subsection{Contributions}
\begin{enumerate}
\item A wallet/EntryPoint contract and bundler for ML-DSA-44 accounts. Key binding, nonce consumption, domain-separated operation hashes and batch-digest recomputation are all enforced on-chain (Section~\ref{sec:system}).
\item A precise statement of the proven relation, its witness and public inputs, the trust model, and a four-layer argument for replay protection (Section~\ref{sec:proof}).
\item An experimental evaluation on a real EVM: repeated runs over nine batch sizes, an ablation separating off-chain verification, Groth16 verification and batching, and an executed native baseline that we compare against the cost model (Section~\ref{sec:eval}).
\item A security and correctness evaluation (\SecCases\ cases), including a negative control that documents the remaining trust assumption (Section~\ref{sec:security}).
\item A quantified off-chain cost analysis, and a reproducible artefact with pinned software (Appendix).
\end{enumerate}

\section{Background and Related Work}\label{sec:related}

\textbf{Account abstraction.} ERC-4337~\cite{eip4337} lets contract wallets define their own validation logic, with user operations collected by bundlers in an alternative mempool. zkSync Era~\cite{zksync2026aa} and Starknet~\cite{starknet2026accounts} provide account abstraction natively at L2. Precompiles for other classical schemes already exist, e.g.\ secp256r1 (RIP-7212~\cite{rip7212}, 3,450 gas; EIP-7951~\cite{eip7951}, 6,900 gas on mainnet).

\textbf{PQ signatures verified in the EVM.} ZKNox's ETHDILITHIUM~\cite{zknox2026ethdilithium} reports 8.1M gas for its NIST-compatible ML-DSA verifier and 4.9M gas for a variant with an EVM-friendly hash. Its experimental packed verifier, and the Fireblocks verifier~\cite{fireblocks2026evmmldsa}, report about 1.2M gas. ETHFALCON~\cite{zknox2026ethfalcon} reports 1.5--3.9M gas for Falcon variants. None of these is audited. Precompile proposals estimate 4,500 gas for ML-DSA-44 (EIP-8051~\cite{eip8051}), 3,000 for Falcon-512 (EIP-8052~\cite{eip8052}), and 6,500 base for pure FIPS~204 ML-DSA-44 (EIP-8355~\cite{eip8355}). An NTT precompile, EIP-7885~\cite{eip7885}, has also been proposed. Schemitt \emph{et al.}~\cite{schemitt2025assessing} measure PQ signing and verification times and simulate their effect on blockchains, but report no EVM gas.

\textbf{zkVMs and SNARK verification.} Groth16~\cite{groth2016size} proofs verify with a constant number of pairings. Under EIP-1108~\cite{eip1108}, a 4-pair check costs 181,000 gas, plus 6,150 gas per public input. RISC Zero~\cite{bruestle2023risczero} and SP1 prove RISC-V execution and wrap their proofs in Groth16 for on-chain verification. SP1 documents about 270k gas for this~\cite{succinct2026sp1prooftypes}. RISC Zero states that its Groth16/BN254 layer depends on a trusted setup and is not quantum-safe~\cite{risczero2026security}. A non-peer-reviewed benchmark reports 5.35\,s to prove ML-DSA-44 verification in RISC Zero on an RTX~4090, without the Groth16 stage~\cite{kaspakii2026pqzkbench}. Jolt~\cite{arun2024jolt} is a lookup-based zkVM. zkLogin~\cite{baldimtsi2024zklogin} authorises transactions with Groth16 proofs about OpenID credentials.

\textbf{PQ signatures elsewhere.} Algorand uses Falcon keys in its state proofs~\cite{algorand2026stateproofs,micali2021compact}, and QRL uses XMSS~\cite{qrl2026docs}. For Ethereum's consensus layer, hash-based multi-signatures aggregated with SNARKs are being designed~\cite{drake2025hashbased,drake2025leansig,leanmultisig2026}. Falcon signatures can be aggregated with LaBRADOR~\cite{aardal2024labrador,sanso2026falconagg}. zkDilithium~\cite{policharla2023pqprivacypass} proves a STARK-friendly Dilithium variant inside a STARK.

Table~\ref{tab:related} places Q-EVM among these approaches. Unlike precompiles it needs no fork. Unlike in-EVM verifiers, its on-chain cost per operation shrinks with batching. Unlike a sound zkVM design, it currently trusts the bundler for signature validity.

\begin{table*}[t]
\centering
\caption{Comparison with related approaches. Gas figures are as reported by each source (unaudited where noted); Q-EVM figures are measured in this work.}\label{tab:related}
\small
\resizebox{\textwidth}{!}{\begin{tabular}{p{3.6cm}p{2.6cm}p{3.6cm}p{1.2cm}p{4.3cm}}
\toprule
Approach & Scheme & On-chain gas per signature & Fork? & Trust / status \\
\midrule
ZKNox ETHDILITHIUM~\cite{zknox2026ethdilithium} & ML-DSA-44 (NIST) & 8.1M (main); 1.2M (exp.\ branch); \emph{executed here:} \NativeExecMain\ / \NativeExec & No & code only; unaudited \\
Fireblocks verifier~\cite{fireblocks2026evmmldsa} & ML-DSA-44 & 1,226,311 + 4.1M once per key & No & code only; unaudited \\
ZKNox ETHFALCON~\cite{zknox2026ethfalcon} & Falcon-512 variants & 1.5M--3.9M & No & code only; unaudited \\
EIP-8051 / EIP-8355~\cite{eip8051,eip8355} & ML-DSA-44 & 4,500 / 6,500 base (estimates) & Yes & clients; Draft \\
EIP-8052~\cite{eip8052} & Falcon-512 & 3,000 (estimate) & Yes & clients; Draft \\
SP1 + Groth16~\cite{succinct2026sp1prooftypes} & any, in guest & $\approx$270k per proof & No & zkVM + trusted setup (BN254, not PQ) \\
Q-EVM, single op (this work) & ML-DSA-44 & \SingleOpGas\ & No & bundler$^\dagger$ + Groth16 setup \\
Q-EVM, batch of \BatchMaxN\ (this work) & ML-DSA-44 & \BatchMaxGas\ per op & No & bundler$^\dagger$ + Groth16 setup \\
\bottomrule
\end{tabular}}\\[2pt]
\footnotesize $^\dagger$The receipt circuit does not verify ML-DSA (Section~\ref{sec:proof}); with the zkVM extension (Section~\ref{sec:future}) the bundler is no longer trusted for validity.
\end{table*}

\section{Design Choices}\label{sec:design}

\textbf{ML-DSA over Falcon.} Falcon-512 signatures are smaller (666 vs.\ 2,420 bytes), but Falcon signing requires floating-point Gaussian sampling. ML-DSA works over $\mathbb{Z}_q[X]/(X^{256}+1)$ with $q=8{,}380{,}417$ using only integer arithmetic, which is simpler to prove in a zkVM or circuit and simpler for hardware signers. We use the final FIPS~204 standard (ML-DSA-44, pure mode, empty context).

\textbf{Forkless deployment.} A precompile would be the cheapest option per signature, but it needs a coordinated hard fork. Q-EVM requires only contracts and off-chain infrastructure.

\textbf{State on-chain, content in the proof.} Replay protection and key binding are enforced by the wallet's storage, not by the proof. They therefore stay correct when the proof system changes, and also against Groth16's proof malleability (Section~\ref{sec:proof}).

\textbf{Batching.} One authorisation per batch amortises the verifier's fixed cost. The per-operation work that remains in the wallet sets the floor that batching can approach.

\section{System}\label{sec:system}

\begin{table}[t]
\centering
\caption{Implementation status of each component.}\label{tab:status}
\small
\begin{tabular}{p{4.1cm}p{3.6cm}}
\toprule
Component & Status \\
\midrule
ML-DSA-44 sign/verify (FIPS 204) & implemented, measured \\
Off-chain verification in bundler & implemented, measured \\
Wallet: key binding, nonces, digest & implemented, measured (revm) \\
Groth16 receipt (binding circuit) & implemented, measured \\
Groth16 verifier contract & implemented, measured (revm) \\
Batching & implemented, measured \\
Native ML-DSA in EVM & \ifexecbaseline executed (third-party verifier) and modelled\else modelled only\fi \\
ML-DSA proven inside a zkVM & \textbf{future work} \\
RISC Zero cycle count, 12.4\,s proving (earlier draft) & not reproducible; not used \\
\bottomrule
\end{tabular}
\end{table}

\subsection{Operation hash}
For an operation by \texttt{sender} with \texttt{nonce} and call data $m$, the user signs
\begin{multline}
h = \mathrm{keccak}(\mathrm{TAG}\,\|\,\mathit{chainid}\,\|\,\mathit{wallet}\,\|\,\mathit{sender}\\
\|\,\mathit{nonce}\,\|\,\mathrm{keccak}(m)\,\|\,\mathrm{keccak}(pk)),
\end{multline}
where each field is one 32-byte word and $\mathrm{TAG}=\mathrm{keccak}(\texttt{"QEVM\_USEROP\_V1"})$. A batch of $N$ operations is identified by $d=\mathrm{keccak}(h_1\|\cdots\|h_N)$.

\subsection{Bundler}
Algorithm~\ref{alg:validate} lists the validation steps. All state the bundler checks is read from the chain, so an operation that has been included is rejected even by a bundler that has lost its mempool.

\begin{algorithm}[t]
\caption{Bundler validation and batching}\label{alg:validate}
\begin{algorithmic}[1]
\Function{Submit}{$op$}
\State \textbf{reject} if payload missing, or $op.chainid$ or $op.wallet$ wrong
\State \textbf{reject} if $\mathit{pk}[op.sender] \neq \mathrm{keccak}(op.pk)$ \Comment{on-chain}
\State \textbf{reject} if $(op.sender, op.nonce)$ seen before
\State \textbf{reject} if $op.nonce \neq \mathit{nonce}[op.sender] + \mathit{pending}(op.sender)$
\State \textbf{reject} if $\neg$\,\textsc{ML-DSA.Verify}$(op.pk, h(op), op.\sigma)$ \Comment{off-chain}
\State add $op$ to mempool
\EndFunction
\Function{Bundle}{}
\State take ops $o_1..o_N$; $d \gets \mathrm{keccak}(h(o_1)\|\cdots\|h(o_N))$
\State $\pi \gets$ \textsc{Groth16.Prove}$(d, p, 1)$
\State send $\pi \| N \| (sender, nonce, \mathrm{keccak}(m), \mathrm{keccak}(pk))_{i}$ to wallet
\EndFunction
\end{algorithmic}
\end{algorithm}

\subsection{Wallet and verifier contracts}
For each operation record, the wallet:
\begin{enumerate}
\item requires $\mathit{pk}[sender]\neq 0$ and $\mathit{pk}[sender]=\mathrm{keccak}(pk)$;
\item requires $\mathit{nonce}[sender]=nonce$ and then increments it;
\item recomputes $h_i$, filling the chain id and its own address from the \texttt{CHAINID} and \texttt{ADDRESS} opcodes.
\end{enumerate}
It then computes $d$, calls the Groth16 verifier with the five public inputs $(d_{hi}, d_{lo}, p_{hi}, p_{lo}, 1)$, and reverts the whole batch, including every nonce update, if any check fails. The verifier computes $vk_x=IC_0+\sum_i x_i IC_i$ with the \texttt{ecMul}/\texttt{ecAdd} precompiles and checks $e(-A,B)\,e(\alpha,\beta)\,e(vk_x,\gamma)\,e(C,\delta)=1$ with the pairing precompile. Both contracts are emitted directly as EVM bytecode. Executing the call data is outside the scope of the measurements, because it costs the same in every configuration we compare.

\section{What Is Proven}\label{sec:proof}

\textbf{Implemented relation.} Over the BN254 scalar field $\mathbb{F}_r$, the public input is $x=(d_{hi}, d_{lo}, p_{hi}, p_{lo}, v)$, where the two limbs of each 32-byte value are its 128-bit halves (an injective encoding, since $2^{128}<r$), $p=\mathrm{keccak}(\texttt{"qevm.mldsa.verify"})$, and $v$ is a validity flag. The witness is $w=(s_1,\dots,s_5)$, and
\begin{equation}
R_{\mathrm{cur}}=\{(x,w): s_i=x_i^2\ (i=1..5),\ v=1\}.
\end{equation}
The circuit has 6 R1CS constraints and 5 public inputs. The setup is circuit-specific and single-party, from a fixed seed (development only).

\textbf{Guarantees.} By knowledge soundness, a valid proof is bound to exactly $(d,p,v{=}1)$: it does not verify for any other digest. This binding comes from the verifying key's $IC$ terms. \textbf{Non-guarantees.} The witness can be computed from $x$, so \emph{anyone holding the (public) proving key can prove any digest}. The proof therefore does not show that any signature was checked. Signature validity rests on the bundler's off-chain FIPS~204 verification. Groth16 proofs are also malleable, so different proof bytes can verify for the same statement.

\textbf{Why the on-chain cost is still representative.} Groth16 verification cost depends on the number of public inputs, not on circuit size. A zkVM receipt wrapped in Groth16 with the same five-input interface costs the same on-chain. Its proving cost, however, is entirely different (Section~\ref{sec:offchain}).

\textbf{Trust model.} Signature validity: the bundler. Key binding, replay protection, batch integrity: the wallet on-chain, no trust needed. Digest binding: Groth16 soundness and the setup. The BN254 layer is classical, so a quantum adversary could forge receipts~\cite{risczero2026security}. With $R_{\mathrm{cur}}$ this changes nothing, because forging is already easy. It does matter for the future zkVM design.

\textbf{Replay protection} has four layers. (1)~\emph{Statement binding}: the proof binds $d$, and $d$ binds every $h_i$, so any edit, reordering, removal, addition or duplication invalidates it. (2)~\emph{Cross-chain}: $h$ contains the chain id, which the wallet takes from \texttt{CHAINID}. (3)~\emph{Cross-contract}: $h$ contains the wallet address, which the wallet takes from \texttt{ADDRESS}. (4)~\emph{Same-operation replay}: the wallet consumes the nonce. Layer~(4) is required because a re-randomised proof is a new, valid proof for the same statement (tested in Section~\ref{sec:security}).

\textbf{Target relation (future work).} $R_{\mathrm{ML\text{-}DSA}}$ takes $x=(d,\mathit{imageID})$ and witness $(pk_i,\sigma_i,h_i)_{i\le N}$, and holds iff $d=\mathrm{keccak}(h_1\|\cdots\|h_N)$ and $\textsc{ML-DSA.Verify}(pk_i,h_i,\sigma_i)=1$ for all $i$. Executed in a zkVM and wrapped in Groth16, it keeps the same on-chain interface and removes the bundler from the trust model.

\section{Experimental Evaluation}\label{sec:eval}

\subsection{Method}
All contracts run on revm~14.0.3 with Cancun gas rules~\cite{revm}, with state committed between transactions. Reported gas is the EVM's \texttt{gas\_used}, including the 21,000 base cost and calldata. Configurations:
\begin{itemize}
\item \textbf{A0}: native ML-DSA in the EVM, one transaction per operation.
\item \textbf{A1}: off-chain verification plus an ECDSA attestation by the bundler, checked with \texttt{ecrecover}. This configuration fully trusts the bundler.
\item \textbf{A2}: off-chain verification plus one Groth16 receipt and one wallet transaction per operation.
\item \textbf{A3}: as A2, but one receipt and one transaction per batch of $N$.
\end{itemize}
A1--A3 run through the same bundler and wallet code and differ only in how the batch is authorised. For each configuration and $N\in\{1,2,4,\dots,256\}$ we discard 3 warm-up runs and keep \Repeats. Every run uses fresh ML-DSA-44 accounts registered with nonce~1. Statistics are the mean, standard deviation, median, p95 and a 95\% Student-$t$ confidence interval. Hardware and software versions are listed in the Appendix.

\subsection{Gas per operation and batching}
\begin{figure}[t]
\centering
\begin{tikzpicture}
\begin{loglogaxis}[width=\columnwidth, height=5.4cm, xlabel={batch size $N$}, ylabel={gas per operation},
  log basis x=2, legend columns=2, legend style={font=\scriptsize, at={(0.5,-0.28)}, anchor=north}, grid=major, ymin=5000, ymax=2e7]
\addplot+[mark=none, dashed, black] table[x=n, y=gas_per_op, col sep=comma] {data/A0-native-model.csv};
\addlegendentry{A0 native (model)}
\ifexecbaseline
\addplot+[mark=none, densely dotted, thick, red] coordinates {(1,\NativeExecRaw) (256,\NativeExecRaw)};
\addlegendentry{A0$'$ native, executed (optimised)}
\addplot+[mark=none, dashdotted, thick, brown] coordinates {(1,\NativeExecMainRaw) (256,\NativeExecMainRaw)};
\addlegendentry{A0$''$ native, executed (main)}
\fi
\addplot+[mark=square*, mark size=1.4pt] table[x=n, y=gas_per_op, col sep=comma] {data/A2-groth16-per-op.csv};
\addlegendentry{A2 Groth16 per op}
\addplot+[mark=*, mark size=1.4pt] table[x=n, y=gas_per_op, col sep=comma] {data/A3-groth16-batch.csv};
\addlegendentry{A3 Groth16 + batching}
\addplot+[mark=triangle*, mark size=1.6pt] table[x=n, y=gas_per_op, col sep=comma] {data/A1-attestation.csv};
\addlegendentry{A1 ECDSA attestation}
\end{loglogaxis}
\end{tikzpicture}
\caption{Measured gas per operation vs.\ batch size (mean of \Repeats\ runs; the 95\% CIs are narrower than the markers).}\label{fig:gas}
\end{figure}

Fig.~\ref{fig:gas} and Table~\ref{tab:sweep} show the batch sweep. A single-operation Q-EVM transaction costs \SingleOpGas\ gas. About 181k of this is the fixed pairing check, and the rest is the five \texttt{ecMul}/\texttt{ecAdd} calls, the calldata, and the wallet's per-operation storage reads and write. With batching, the fixed cost is shared: gas per operation falls to \BatchEightGas\ at $N=8$ and \BatchMaxGas\ at $N=\BatchMaxN$. It approaches the per-operation floor set by the wallet's two cold storage reads, one storage write, operation-hash recomputation and 128 bytes of calldata. The attestation wallet (A1) reaches \AttestMaxGas\ gas per operation at the same $N$.

\begin{table}[t]
\centering
\caption{A3 (Q-EVM) gas per operation vs.\ $N$; reduction relative to the native baseline\ifexecbaseline\ (model and executed)\fi.}\label{tab:sweep}
\scriptsize
\resizebox{\columnwidth}{!}{\begin{tabular}{rrrrrrr}
\toprule
$N$ & \multicolumn{2}{c}{A3 gas/op} & A1 gas/op & A3 vs model & A3 vs exec. & e2e/op (ms) \\
 & mean & 95\% CI & & & & \\
\midrule
1 & 251{,}302 & $\pm$4 & 34{,}953 & 91.0\% & 79.9\% & 38.69 \\
2 & 130{,}204 & $\pm$3 & 22{,}035 & 95.4\% & 89.6\% & 20.61 \\
4 & 69{,}660 & $\pm$2 & 15{,}574 & 97.5\% & 94.4\% & 9.89 \\
8 & 39{,}387 & $\pm$1 & 12{,}344 & 98.6\% & 96.8\% & 5.00 \\
16 & 24{,}249 & $\pm$1 & 10{,}728 & 99.1\% & 98.1\% & 2.66 \\
32 & 16{,}681 & $\pm$1 & 9{,}921 & 99.4\% & 98.7\% & 1.34 \\
64 & 12{,}898 & $\pm$0 & 9{,}517 & 99.5\% & 99.0\% & 0.71 \\
128 & 11{,}005 & $\pm$0 & 9{,}315 & 99.6\% & 99.1\% & 0.42 \\
256 & 10{,}060 & $\pm$0 & 9{,}215 & 99.6\% & 99.2\% & 0.27 \\
\bottomrule
\end{tabular}
}
\end{table}

\subsection{Ablation}
Table~\ref{tab:ablation} separates the three mechanisms:
\begin{itemize}
\item \emph{Off-chain verification} (A0$\to$A1) removes the in-EVM lattice arithmetic and hashing, which dominates every other cost. It needs only a trusted bundler.
\item \emph{Groth16 verification} (A1$\to$A2, $N=1$) costs \GrothOverhead\ gas more per operation than an attestation (\PerOpGas\ vs.\ \AttestOneGas). In return it replaces a bundler signature with a succinct proof whose interface is identical to that of a sound zkVM receipt.
\item \emph{Batching} (A2$\to$A3) removes most of the Groth16 overhead as $N$ grows.
\end{itemize}
With the current binding circuit, A2/A3 and A1 have the same trust assumption (the bundler). The proof only begins to pay for its gas once the circuit is sound (Section~\ref{sec:future}).

\begin{table}[t]
\centering
\caption{Ablation: gas per operation (mean of \Repeats\ runs).}\label{tab:ablation}
\scriptsize
\resizebox{\columnwidth}{!}{\begin{tabular}{lrrrrl}
\toprule
Configuration & $N=1$ & $N=8$ & $N=32$ & $N=256$ & Trust \\
\midrule
A0$'$ native, executed (optimised) & 1{,}248{,}996 & 1{,}248{,}996 & 1{,}248{,}996 & 1{,}248{,}996 & none \\
A0$''$ native, executed (main) & 8{,}168{,}152 & 8{,}168{,}152 & 8{,}168{,}152 & 8{,}168{,}152 & none \\
A0 native ML-DSA (model) & 2{,}802{,}020 & 2{,}802{,}047 & 2{,}802{,}048 & 2{,}802{,}032 & none \\
A1 off-chain + ECDSA attest. & 34{,}953 & 12{,}344 & 9{,}921 & 9{,}215 & bundler key \\
A2 + Groth16, per op & 251{,}297 & 251{,}302 & 251{,}299 & 251{,}300 & bundler$^\dagger$ \\
A3 + Groth16 + batching & 251{,}302 & 39{,}387 & 16{,}681 & 10{,}060 & bundler$^\dagger$ \\
\bottomrule
\end{tabular}
}\\[2pt]
\footnotesize $^\dagger$receipt circuit does not verify ML-DSA.
\end{table}

\subsection{Native baseline: modelled vs.\ executed}
The earlier draft compared measured Q-EVM gas against a \emph{modelled} native cost (about 2.8M gas, from an opcode model calibrated to a published estimate), so the two sides of the comparison were obtained by different methods. We now report both a model and executed baselines. The model is computed on the same keys and signatures as the experiments (A0, \NativeModel\ gas).
\ifexecbaseline
In addition, we \emph{executed} two versions of the open-source ZKNox NIST ML-DSA-44 verifier~\cite{zknox2026ethdilithium} on the \emph{same} revm chain and hard fork, with the same 100 FIPS~204 keys, messages and signatures (Table~\ref{tab:native}):
\begin{itemize}
\item the optimised experimental verifier (branch \texttt{exp/packed-verifier}, with its Keccak-f[1600] helper contract): \NativeExec\ gas per \texttt{verify} transaction;
\item the main-branch verifier: \NativeExecMain\ gas.
\end{itemize}
Both accepted every valid signature and rejected every tampered one. Both also need a one-off \texttt{setKey} of about \SetKeyGas\ gas per account, to store the expanded public key. The model lies between the two executed verifiers ($\ExecVsModel\%$ and $\MainVsModel\%$ off), which confirms that a modelled baseline cannot replace a measured one. We use the cheapest executed verifier as the primary baseline, because it is the least favourable comparison for Q-EVM. Against it, Q-EVM saves \RedExecOne\% for a single operation and \RedExecMax\% per operation at $N=\BatchMaxN$. Against the main-branch verifier, it saves \RedMainOne\% and \RedMainMax\%. The earlier draft's single figure of 91.4\% thus depended entirely on the choice of baseline.
\begin{table}[t]
\centering
\caption{Executed native ML-DSA-44 baselines on the same revm chain (100 signatures, gas includes the 21,000 base cost and calldata).}\label{tab:native}
\scriptsize
\resizebox{\columnwidth}{!}{\begin{tabular}{lrr}
\toprule
Quantity & optimised$^a$ & main branch$^b$ \\
\midrule
\texttt{verify} tx gas, mean ($n=100$) & 1{,}248{,}996 & 8{,}168{,}152 \\
\quad min & 1{,}244{,}485 & 8{,}079{,}771 \\
\quad max & 1{,}252{,}253 & 8{,}322{,}498 \\
One-off \texttt{setKey} per account & 4{,}837{,}624 & 4{,}837{,}676 \\
Opcode cost model, same inputs & \multicolumn{2}{c}{2{,}801{,}997} \\
Model error vs.\ executed & +124.3\% & -65.7\% \\
\bottomrule
\end{tabular}
}\\[2pt]
\footnotesize $^a$ZKNox \texttt{exp/packed-verifier} @\,\texttt{f85ca80}, solc 0.8.30 via-IR, 10$^6$ runs. $^b$ZKNox main @\,\texttt{f9f4fe1}, solc 0.8.30, $10^4$ runs. Both compiled for Cancun.
\end{table}
\else
An executed baseline (a third-party Solidity verifier run on the same chain) could not be completed within our time box, so the native figure remains a \emph{model}. Published Solidity verifiers report 1.2M--8.1M gas~\cite{fireblocks2026evmmldsa,zknox2026ethdilithium}. Even against the cheapest of these, the batched cost of Q-EVM is more than two orders of magnitude lower.
\fi

\subsection{Time, latency and memory}
Table~\ref{tab:time} reports wall-clock times. Off-chain FIPS~204 verification takes \MldsaVerifyMsPerOp\,ms per signature. Proving the binding circuit takes \ProveMsOne\,ms, independent of $N$, because the circuit is fixed. End-to-end bundler time, from the first submission until the batch is committed on-chain, is \EtoEMsOne\,ms for $N=1$ and \EtoEMsMax\,ms for $N=\BatchMaxN$ (\EtoEPerOpMax\,ms per operation). Peak heap stays below \HeapMax\,MiB. These times describe the \emph{current} system and are not the cost of proving ML-DSA (Section~\ref{sec:offchain}).

\begin{table*}[t]
\centering
\caption{Off-chain time and memory (mean of \Repeats\ runs; laptop, see Appendix). ``Authorise'' is Groth16 proving of the binding circuit (A2/A3) or ECDSA signing (A1).}\label{tab:time}
\small
\begin{tabular}{lrrrrrr}
\toprule
Config & $N$ & ML-DSA verify & Authorise & EVM exec & End-to-end & Peak heap \\
 & & (ms, all $N$) & (ms) & (ms) & (ms, mean $\pm$ sd) & (MiB) \\
\midrule
A1 & 1 & 0.03 & 0.05 & 0.22 & 0.6 $\pm$ 0.0 & 0.4 \\
A1 & 8 & 0.23 & 0.06 & 0.24 & 1.4 $\pm$ 0.1 & 0.7 \\
A1 & 32 & 0.95 & 0.06 & 0.31 & 4.1 $\pm$ 0.5 & 1.6 \\
A1 & 256 & 7.37 & 0.07 & 0.95 & 27.8 $\pm$ 1.2 & 10.1 \\
\midrule
A2 & 1 & 0.03 & 25.55 & 6.81 & 46.4 $\pm$ 2.2 & 0.6 \\
A2 & 8 & 0.38 & 194.37 & 55.62 & 363.2 $\pm$ 11.5 & 0.9 \\
A2 & 32 & 1.62 & 796.99 & 230.80 & 1497.0 $\pm$ 53.7 & 1.8 \\
A2 & 256 & 12.86 & 6595.27 & 1679.38 & 12189.1 $\pm$ 191.1 & 10.5 \\
\midrule
A3 & 1 & 0.03 & 24.13 & 6.90 & 38.7 $\pm$ 2.0 & 0.6 \\
A3 & 8 & 0.24 & 24.51 & 7.02 & 40.0 $\pm$ 2.0 & 0.9 \\
A3 & 32 & 0.97 & 24.55 & 7.15 & 42.7 $\pm$ 1.8 & 1.8 \\
A3 & 256 & 7.79 & 23.06 & 7.49 & 68.0 $\pm$ 16.3 & 10.3 \\
\bottomrule
\end{tabular}

\end{table*}

\subsection{Cost shifted off-chain}\label{sec:offchain}
A sound receipt requires proving ML-DSA verification inside a zkVM. We did not run that, because it needs a Linux GPU host (Section~\ref{sec:future}). Instead we treat the per-signature proving time $t_{\mathrm{sig}}$ as a parameter. The earlier draft's projection of 12.4\,s is one value of it, as is the third-party measurement of 5.35\,s without the Groth16 wrap~\cite{kaspakii2026pqzkbench}. At a prover price $c$ (USD/s; an AWS g5.xlarge costs 1.006\,USD/h~\cite{aws2026g5pricing}), the proof costs $c\,t_{\mathrm{sig}}$ per operation. At gas price $g$ and ether price $E$, it saves $(G_{\mathrm{native}}-G_{\mathrm{Q\text{-}EVM}}(N))\,g\,E$. Batching lowers $G_{\mathrm{Q\text{-}EVM}}(N)$ but not $t_{\mathrm{sig}}$. On one prover, a batch therefore waits at least $N\,t_{\mathrm{sig}}$ before inclusion: batching trades on-chain gas for latency, and only parallel provers reduce wall-clock time. Table~\ref{tab:offchain} evaluates this against the executed optimised baseline at $N=\BatchMaxN$, with an ether price of \EthUsd\,USD (an assumption). At the projected $t_{\mathrm{sig}}=12.4$\,s, the proof costs \ProverUsdPerOpProj\,USD per operation, while the gas saved is worth \SavedUsdOneGwei\,USD at 1\,gwei. Proving would only cost as much as it saves at $t_{\mathrm{sig}}\approx$ \BreakEvenOneGwei\,s (1\,gwei) or \BreakEvenTenthGwei\,s (0.1\,gwei). \emph{Monetary} cost is therefore not the constraint. \emph{Latency} is: proving must finish before inclusion, so a batch of $N$ waits at least $N\,t_{\mathrm{sig}}/P$ on $P$ parallel provers. This confines a sound Q-EVM to non-interactive, high-value or batched settlement unless proving becomes much faster.

\begin{table}[t]
\centering
\caption{Net saving per operation (USD) of a sound zkVM proof at $N=\BatchMaxN$ vs.\ the executed optimised native verifier; prover at \ProverUsdHour\,USD/h (AWS g5.xlarge on-demand~\cite{aws2026g5pricing}). $t_{sig}$ is a parameter; 12.4\,s is the earlier draft's projection.}\label{tab:offchain}
\scriptsize
\begin{tabular}{rrrr}
\toprule
$t_{sig}$ (s), prover USD/op & 0.1 gwei & 1 gwei & 10 gwei \\
\midrule
1, 0.0003 & 0.371 & 3.717 & 37.168 \\
5, 0.0014 & 0.370 & 3.715 & 37.167 \\
12.4, 0.0035 & 0.368 & 3.713 & 37.165 \\
30, 0.0084 & 0.363 & 3.708 & 37.160 \\
60, 0.0168 & 0.355 & 3.700 & 37.151 \\
\bottomrule
\end{tabular}

\end{table}

\section{Security and Correctness Evaluation}\label{sec:security}
We ran \SecCases\ cases through the real pipeline: bundler validation, off-chain FIPS~204 verification, and the wallet on revm (Table~\ref{tab:sec}). The cases cover:
\begin{itemize}
\item valid batches of 1, 8 and 64 operations;
\item invalid signatures: bit flips in $\tilde c$, $\mathbf z$ and $\mathbf h$, malformed hints, truncated, extended or all-zero signatures, another key, another message;
\item operations modified after signing (call data, nonce, chain id, wallet, sender, key);
\item key binding;
\item replays before bundling, after inclusion, to a fresh bundler, of the identical wallet transaction, and with a \emph{re-randomised} Groth16 proof;
\item altered batches (reordered, dropped, added or duplicated operations, wrong $N$, truncated);
\item invalid proofs (zero, random, points off the curve, another setup, another program, another chain or wallet).
\end{itemize}
\SecPass/\SecCases\ cases behave as expected. Unit tests additionally cover a $\mathbb{G}_2$ point outside the subgroup, public inputs $\ge r$, and full rollback of a failed batch.

The \emph{negative control} is an unsigned operation with a proof made directly from the public proving key, and the wallet \emph{accepts} it. It documents the remaining trust assumption of Section~\ref{sec:proof}, and must flip to \emph{rejected} once ML-DSA is proven inside a zkVM. A rejected proof makes the precompile call consume all the gas it was given (about 29M in our runs), which the submitter pays. Bundlers must therefore simulate before submitting, as ERC-4337 bundlers already do.

\begin{table}[t]
\centering
\caption{Security and correctness evaluation (details: \texttt{results/security.md}).}\label{tab:sec}
\scriptsize
\resizebox{\columnwidth}{!}{\begin{tabular}{p{2.3cm}rp{3.6cm}l}
\toprule
Category & Cases & Decided by & As expected \\
\midrule
valid & 3 & bundler+wallet & 3/3 \\
invalid signature & 9 & bundler & 9/9 \\
modified op hash & 6 & bundler & 6/6 \\
key binding & 2 & bundler & 2/2 \\
replay & 8 & bundler, bundler (reads on-chain nonce), wallet & 8/8 \\
altered batch & 7 & wallet & 7/7 \\
invalid proof & 9 & wallet & 9/9 \\
soundness (negative control) & 1 & wallet & accepted (gap C1) \\
\midrule
Total & 45 & & 45/45 \\
\bottomrule
\end{tabular}
}
\end{table}

\section{Limitations}
\begin{itemize}
\item Signature validity is trusted to the bundler (Section~\ref{sec:proof}).
\item The Groth16 layer is classical (BN254) and uses a development setup.
\item Gas is measured on revm with Cancun rules. Pectra's calldata floor (EIP-7623~\cite{eip7623}) is not applied. A sound design must also respect Fusaka's per-transaction cap of $2^{24}$ gas (EIP-7825~\cite{eip7825}), which already limits in-EVM verifiers to about two signatures per transaction.
\item Timings come from a laptop (AMD Ryzen~5 5500U) inside a container.
\item The wallet omits call-data execution and fee payment, which are equal across configurations.
\item No testnet deployment was made.
\end{itemize}

\section{Future Work}\label{sec:future}
\begin{enumerate}
\item \textbf{Sound proof (FW-1).} Run FIPS~204 verification in a RISC~Zero or SP1 guest on a Linux x86-64 host with at least 32\,GB of RAM and optionally a GPU. Measure cycles, STARK proving time, Groth16 wrap time and peak memory for $N\in\{1,\dots,16\}$. Verify the result with the production verifier on the same chain, and add ablation rows A4 (zkVM, per operation) and A5 (zkVM, batched). The negative control must then be rejected.
\item Re-run all measurements with Prague/Pectra rules.
\item A custom ML-DSA circuit.
\item Recursive aggregation of batch receipts.
\item A multi-party trusted setup, or a transparent proof system.
\item A testnet deployment.
\item Decentralised proving.
\end{enumerate}

\section{Conclusion}
Moving ML-DSA verification off-chain and authorising batches with a succinct receipt reduces the on-chain cost of a post-quantum account from \NativeExec--\NativeExecMain\ gas per signature (executed Solidity verifiers) to \SingleOpGas\ gas for one operation and \BatchMaxGas\ gas per operation at $N=\BatchMaxN$. The wallet enforces key binding and replay protection on-chain. The price is off-chain infrastructure, latency, and, in the current implementation, trust in the bundler for signature validity. The key next step is to remove that trust with a zkVM proof of ML-DSA verification, at a proving cost that our analysis shows to be economically acceptable for batched, high-value settlement.

\appendix[Reproducibility]
The code, results and scripts are available at~\cite{qevmrepo}.
\begin{itemize}
\item \emph{Hardware:} AMD Ryzen~5 5500U (6 cores, 12 threads), 8\,GB RAM, no discrete GPU. Windows~11 host; runs inside a Debian~12 container (Docker 29.8.1), about 1.9\,GB of RAM visible.
\item \emph{Software:} Rust 1.100.0; revm 14.0.3 (Cancun); pqcrypto-mldsa 0.1.2~\cite{pqcryptomldsa}; ark-groth16/ark-bn254 0.4.0~\cite{arkworks}; k256 0.13.4; all versions pinned in \texttt{Cargo.lock}. Baseline: ZKNox ETHDILITHIUM at commit \texttt{f9f4fe1}, compiled with solc 0.8.30 (\texttt{--optimize-runs 10000 --evm-version cancun}).
\item \emph{Commands:} \texttt{qevm-cli security-report}, \texttt{qevm-cli experiment --repeats 30}, \texttt{scripts/native-baseline.sh}; tables are generated by \texttt{scripts/make\_paper\_tables.py}.
\item \emph{Gas schedule:} EIP-2028 calldata, EIP-1108 BN254 precompiles, EIP-2929 storage access.
\item \emph{Mainnet context:} block gas limit 60,000,000 as of 2026-11-15.
\end{itemize}

\bibliographystyle{IEEEtran}
\bibliography{refs}

\end{document}
